\documentclass{misisarticle}
\grnti{28.23.37}
\asjc{1702, 1712}
\oecd{1.02}
\title{Переусложнение мультиагентных схем генерации кода: зависимость эффективности от сложности задачи и силы базовой модели}
\misisauthor{Павел Валерьевич Дурынин}{кафедра инженерной кибернетики НИТУ «МИСиС»}{Москва, Россия}{}
\abstract{Рассматривается задача выбора между одноагентной и мультиагентной генерацией программного кода. Проверяется гипотеза о том, что дополнительные роли и обмен сообщениями могут не давать выигрыша на простых задачах и сильных базовых моделях при увеличении затрат. Предлагается протокол сравнения прямой генерации, одноагентной самокоррекции, ролевого взаимодействия и ансамбля кандидатов на моделях Qwen2.5-Coder разного размера и задачах EvalPlus и LiveCodeBench. Эффективность оценивается по доле корректных итоговых программ, расходу токенов и времени выполнения. Предусмотрены парные статистические проверки и анализ траекторий по таксономии отказов MAST. Планируемый результат исследования представляет собой карту применимости схем в координатах «сложность задачи~--- сила модели». В статье представлены постановка задачи и методика исследования; экспериментальные результаты не приводятся.}
\keywords{генерация программного кода; мультиагентные системы; одноагентная самокоррекция; переусложнение; сложность задачи; масштаб языковой модели}
\begin{document}
\maketitle
\section{Введение}
Генерация программного кода по текстовому условию допускает несколько способов организации работы языковой модели. В простейшем случае модель сразу возвращает программу. Более сложные схемы добавляют планирование, создание тестов, исправление ошибок и выбор из нескольких решений. Отдельные вызовы могут получать разные роли и независимые контексты, образуя мультиагентную систему. При этом увеличение числа вызовов одновременно меняет организацию решения и объём затраченных ресурсов.

Для задач программирования разработаны специализированные мультиагентные подходы AgentCoder и MapCoder~\cite{agentcoder,mapcoder}. Однако результат конкретной схемы на выбранном наборе задач не означает, что она предпочтительна при любой сложности задания и для любой базовой модели. Исследование Gao и соавторов указывает, что преимущество мультиагентных систем может сокращаться с ростом возможностей модели~\cite{hybrid}. Возникает вопрос: в каких условиях координация помогает получить корректную программу, а в каких расходует дополнительные ресурсы без полезного изменения результата?

В настоящей работе переусложнение понимается как наблюдаемое свойство сравнения схем, а не как заранее установленный недостаток мультиагентности. Различаются два случая: подтверждённое ухудшение корректности и рост затрат при практически эквивалентной корректности. Первый характеризует потерю качества, второй~--- невыгодное использование ресурсов. Они требуют разных критериев.

Цель исследования~--- определить области применимости одноагентных и мультиагентных схем генерации Python-кода в зависимости от сложности задач и возможностей модели. Планируемый вклад состоит в едином протоколе ресурсного сравнения, карте применимости и разборе механизмов ошибок по сохранённым траекториям. Методика ограничена алгоритмическими задачами с автоматической проверкой. Перенос выводов на разработку больших репозиториев требует отдельного исследования.

\section{Обзор работ}
\subsection{Одноагентная генерация и исправление}
Прямая генерация служит исходной точкой: одному условию соответствует одна возвращаемая программа. Такой режим позволяет оценить возможности самой модели до добавления оркестрации. Семейство Qwen2.5-Coder содержит модели нескольких размеров и предназначено для задач, связанных с кодом~\cite{qwen}. Использование одного семейства уменьшает различия между сравниваемыми моделями, хотя размер параметров сам по себе не является точной мерой способности решать задачи.

Self-Debugging использует объяснение программы и обратную связь о выполнении для исправления исходного решения~\cite{selfdebug}. Повторное рассуждение и отладка доступны одной модели без выделения независимых агентов. Поэтому сравнения только с одним прямым вызовом недостаточно: выигрыш ролевой схемы может объясняться дополнительными попытками исправления. Одноагентная самокоррекция выступает самостоятельной базовой схемой.

\subsection{Ролевые схемы и ансамбли}
В AgentCoder выделены программист, проектировщик тестов и исполнитель тестов; результаты выполнения используются для последующего исправления программы~\cite{agentcoder}. Разделение ответственности позволяет отделить создание решения от подготовки проверок. Вместе с тем тесты, созданные моделью, могут быть неполными или неверными, поэтому их нельзя использовать как окончательный эталон корректности.

MapCoder организует работу вокруг воспоминания подходящих примеров, планирования, генерации и отладки~\cite{mapcoder}. Таким образом, термин «мультиагентная генерация» объединяет разные протоколы. Число ролей, доступная обратная связь и правило остановки могут влиять на результат независимо друг от друга. В данной работе используется явно заданная ролевая схема, основанная на этих принципах; она не объявляется точным воспроизведением исходных систем.

Работа More Agents Is All You Need рассматривает независимое семплирование и голосование, а также связь выигрыша со сложностью задачи~\cite{forest}. Для кода буквальное совпадение ответов не подходит: разные программы могут реализовывать один алгоритм. Поэтому ансамбль определяется через поведение кандидатов на доступных проверочных входах и анализируется отдельно от систем с обменом сообщениями.

\subsection{Границы полезности и отказы}
Gao и соавторы сравнивают одноагентные и мультиагентные системы и предлагают гибридное распределение запросов~\cite{hybrid}. Их результаты мотивируют проверку зависимости выигрыша от силы модели, однако не заменяют исследование выбранных схем генерации Python-кода при общей сетке бюджетов.

Cemri и соавторы предлагают таксономию MAST, объединяющую ошибки проектирования системы, несогласованность между агентами и ошибки проверки и завершения~\cite{mast}. Она используется для описания наблюдаемых сбоев. Одна только метка MAST не доказывает, что ошибка вызвана именно мультиагентностью: необходимо сопоставить траектории и провести контролируемое изменение схемы.

Исследовательский вопрос состоит в определении условий полезности архитектуры. Предлагаемый протокол связывает сложность задачи, возможности модели, итоговую корректность и полные затраты на получение одного ответа.

\section{Постановка задачи}
Пусть $D_d$~--- множество задач уровня сложности $d$, а $m$~--- фиксированная версия модели. Задача $x_i$ содержит условие и открытые примеры; закрытые проверки $H_i$ доступны только оценщику. Схема $s$ при ресурсном ограничении $B$ и случайном состоянии $r$ возвращает одну итоговую программу $z_{i,r}^{s,m,B}$.

Вводится индикатор успешного решения
\begin{equation}
 y_{i,r}^{s,m,B}=\mathbf{1}\{z_{i,r}^{s,m,B}\text{ проходит }H_i\}.
\end{equation}
Синтаксическая ошибка, превышение лимита выполнения или отсутствие итоговой программы считаются неуспехом. Ошибка инфраструктуры фиксируется отдельно и обрабатывается по заранее заданному правилу, одинаковому для всех схем.

Для каждой ячейки сравнения оценивается
\begin{equation}
 P_s(m,d,B)=\frac{1}{|D_d|R}\sum_{i\in D_d}\sum_{r=1}^{R}y_{i,r}^{s,m,B},
\end{equation}
где $R$~--- число независимых повторов. Это оценка вероятности успеха одной полной попытки схемы, возвращающей один ответ (pass@1 на уровне схемы). Внутренние кандидаты ансамбля не считаются отдельными ответами. Выбор лучшего кандидата по закрытым тестам запрещён.

Затраты $C_s$ включают все входные и выходные токены всех вызовов, в том числе повторно переданный контекст, сообщения ролей, отладку и выбор ответа. Дополнительно сохраняются время $T_s$, число вызовов и число запусков тестов. Для выбранной одноагентной базы $b$ вычисляются
\begin{equation}
 \Delta P=P_s-P_b,\qquad \Delta C=C_s-C_b.
\end{equation}
Ухудшение качества фиксируется, если верхняя граница доверительного интервала для $\Delta P$ меньше нуля. Практическая эквивалентность проверяется относительно заранее выбранного допуска $\varepsilon$: весь интервал должен лежать внутри $[-\varepsilon,\varepsilon]$. Эквивалентность при подтверждённом росте затрат трактуется как ресурсное переусложнение. Отсутствие значимого выигрыша само по себе не означает эквивалентности.

Задача исследования~--- оценить различия для каждой комбинации модели, сложности и бюджета и выделить области выигрыша, ухудшения, эквивалентности и недостаточной определённости. Проверяется гипотеза, что выигрыш ролевой схемы уменьшается на более простых задачах и при более сильной базовой модели. Гипотеза допускает опровержение.

\section{Описание методов}
\subsection{Сравниваемые схемы}
\textit{Прямая генерация.} Модель получает условие, формат ответа и открытые примеры, затем возвращает одну программу. Это минимальная схема без дополнительных циклов отладки.

\textit{Одноагентная самокоррекция.} Та же модель создаёт программу, получает результат доступных тестов и исправляет решение в общем контексте. Число итераций ограничено. Схема останавливается при выполнении правила успеха на доступных проверках либо при исчерпании лимита. Закрытые тесты на остановку не влияют.

\textit{Ролевая схема.} Независимые контексты одной модели используются для планировщика, программиста и проектировщика тестов. Планировщик предлагает алгоритм, программист пишет код, проектировщик готовит проверки по условию. Детерминированный исполнитель запускает проверки и передаёт отчёт программисту. Исполнитель является инструментом, а не дополнительной языковой моделью. Число циклов, порядок сообщений и состав контекста фиксируются до основного сравнения.

\textit{Ансамбль с голосованием.} Независимые контексты генерируют кандидатов без обмена сообщениями. После отсева программ, не прошедших открытые примеры, кандидаты группируются по векторам выходов на дополнительных допустимых входах. Выбирается первый сгенерированный представитель наибольшей группы; при равенстве размеров групп также используется порядок генерации. Если подходящих кандидатов нет, попытка считается неуспешной. Задачи с несколькими допустимыми ответами исключаются из общего сравнения всех схем по заранее заданному признаку, поскольку совпадение выходов для них не является подходящим правилом голосования.

Исходные открытые проверки и дополнительные входы одинаковы для всех итеративных схем. Если они подготовлены отдельным вызовом модели, его токены входят в затраты каждой использующей их схемы. Входы без известного правильного ответа применяются только для группировки поведения. Сгенерированные утверждения с ожидаемым ответом считаются вспомогательной обратной связью и сохраняются для анализа ошибок. Только закрытые проверки определяют итоговую корректность.

\subsection{Модели и сложность задач}
Планируется использовать инструктивные версии Qwen2.5-Coder размером 1.5B, 3B, 7B, 14B и 32B~\cite{qwen}. Внутри прогона все языковые роли используют один контрольный снимок модели. Фиксируются версия, точность весов, шаблон диалога, параметры семплирования и ограничения контекста. Варианты с разной квантизацией не объединяются в одну серию без отдельной проверки.

Размер модели является управляемым фактором. Её фактическая сила дополнительно характеризуется качеством прямой генерации на отдельной калибровочной выборке, не участвующей в итоговой оценке. Строгая монотонность качества по числу параметров не предполагается.

Основную стратификацию дают метки easy, medium и hard в выбранном выпуске LiveCodeBench~\cite{livecodebench}. HumanEval+ и MBPP+ из EvalPlus используются как дополнительные наборы функциональных задач~\cite{evalplus,evalplusrepo}. Принадлежность к бенчмарку не считается доказательством низкой сложности каждой задачи. Результаты наборов приводятся отдельно, чтобы различия формата не подменяли эффект сложности.

Версии данных и правила исключения фиксируются до запусков. Для LiveCodeBench выбирается временной интервал после опубликованной границы обучающих данных модели. Если граница не подтверждена, временная фильтрация рассматривается только как снижение риска утечки. HumanEval+ и MBPP+ сохраняют риск знакомства модели с исходными задачами.

\subsection{Контроль ресурсов и проверка программ}
Сравнение выполняется в двух режимах. В первом каждая схема работает со своим фиксированным протоколом, а полные затраты измеряются после завершения. Во втором всем схемам задаются одинаковый предел суммарных токенов и одинаковые ограничения на выполнение тестов. Сетка бюджетов определяется на калибровочной выборке. Прямой вызов может не израсходовать разрешённый бюджет целиком, поэтому отдельно сравнивается самокоррекция, способная использовать дополнительные вызовы.

Для всех схем одинаковы условие, язык Python, исходные открытые данные и ограничения исполнения. Число кандидатов и параметры выбираются до проверки на закрытых тестах. Лимит применяется ко всей схеме, а не к каждому агенту отдельно. Ответ, не завершённый к исчерпанию лимита, учитывается как неуспех.

Код выполняется в изолированной среде без сети и секретов, с ограничениями времени, памяти и числа процессов. Закрытые проверки запускаются после фиксации итогового ответа. Ни их содержимое, ни результаты не возвращаются агентам в рамках попытки. Сохраняются код, доступные проверки, все сообщения, причины остановки и расход ресурсов.

\subsection{Статистическое сравнение}
Схемы проверяются на одних задачах и общей сетке случайных состояний. Для $\Delta P$ и $\Delta C$ строятся парные 95-процентные бутстреп-интервалы с повторной выборкой задач целиком: все повторы одной задачи остаются вместе. Такой способ учитывает зависимость повторных решений одного условия.

Для заранее выбранного основного повтора применяется точный критерий Мак-Немара по несовпадающим исходам двух схем на каждой задаче. Остальные повторы служат для проверки устойчивости; они не считаются независимыми дополнительными задачами. При множественных сравнениях внутри заранее определённого семейства гипотез используется поправка Холма. Основная одноагентная база и допуск $\varepsilon$ выбираются до основного анализа.

Карта применимости строится отдельно для каждого бюджета. В ячейке указываются размер выборки, разница корректности с интервалом и разница затрат. Широкие интервалы означают недостаток данных. Дополнительно рассматривается граница Парето: схема доминируется, если другая даёт не меньшее качество при не больших затратах и хотя бы одно преимущество строгое.

\subsection{Анализ ошибок и воспроизводимость}
Неуспешные ролевые траектории размечаются по определениям MAST~\cite{mast}. Для метки сохраняется фрагмент переписки или кода, подтверждающий решение разметчика. Допускаются несколько меток и неопределённый случай. Часть траекторий независимо размечают два человека; публикуются согласованность и порядок разрешения расхождений. Доли ошибок рассчитываются внутри групп с явным указанием знаменателя.

Чтобы отделить ошибку модели от последствий координации, рассматриваются парные случаи: базовая схема решила задачу, а ролевая потеряла правильный алгоритм или заменила корректный код ошибочным. Для проверки механизма планируется абляция с удалением одного этапа при сохранении бюджета и исходных доступных данных. Такой анализ уточняет интерпретацию карты, но не заменяет оценку корректности.

Для воспроизводимости сохраняются версии моделей и данных, запросы, лимиты, случайные состояния, окружение и траектории. Численные результаты и рекомендации по выбору схем будут получены после выполнения описанного протокола.

\begin{thebibliography}{99}
\bibitem{agentcoder} Huang D., Zhang J. M., Luck M., Bu Q., Qing Y., Cui H. AgentCoder: Multi-Agent-based Code Generation with Iterative Testing and Optimisation. 2023. arXiv:2312.13010. \url{https://arxiv.org/abs/2312.13010}.
\bibitem{mapcoder} Islam M. A., Ali M. E., Parvez M. R. MapCoder: Multi-Agent Code Generation for Competitive Problem Solving // Proceedings of ACL. 2024. \url{https://aclanthology.org/2024.acl-long.269/}.
\bibitem{hybrid} Gao M., Li Y., Liu B., Yu Y., Wang P., Lin C.-Y., Lai F. Single-agent or Multi-agent Systems? Why Not Both? 2025. arXiv:2505.18286. \url{https://arxiv.org/abs/2505.18286}.
\bibitem{qwen} Hui B., Yang J., Cui Z. et al. Qwen2.5-Coder Technical Report. 2024. arXiv:2409.12186. \url{https://arxiv.org/abs/2409.12186}.
\bibitem{selfdebug} Chen X., Lin M., Schärli N., Zhou D. Teaching Large Language Models to Self-Debug. 2023. arXiv:2304.05128. \url{https://arxiv.org/abs/2304.05128}.
\bibitem{forest} Li J., Zhang Q., Yu Y., Fu Q., Ye D. More Agents Is All You Need // Transactions on Machine Learning Research. 2024. arXiv:2402.05120. \url{https://arxiv.org/abs/2402.05120}.
\bibitem{mast} Cemri M., Pan M. Z., Yang S. et al. Why Do Multi-Agent LLM Systems Fail? 2025. arXiv:2503.13657. \url{https://arxiv.org/abs/2503.13657}.
\bibitem{livecodebench} Jain N., Han K., Gu A. et al. LiveCodeBench: Holistic and Contamination Free Evaluation of Large Language Models for Code. 2024. arXiv:2403.07974. \url{https://arxiv.org/abs/2403.07974}.
\bibitem{evalplus} Liu J., Xia C. S., Wang Y., Zhang L. Is Your Code Generated by ChatGPT Really Correct? Rigorous Evaluation of Large Language Models for Code Generation. 2023. arXiv:2305.01210. \url{https://arxiv.org/abs/2305.01210}.
\bibitem{evalplusrepo} EvalPlus. Наборы HumanEval+ и MBPP+, средства проверки программ. \url{https://github.com/evalplus/evalplus}. Дата обращения: 06.11.2026.
\end{thebibliography}
\end{document}
